\documentclass[11pt]{article}
\usepackage[margin=25mm]{geometry}
\usepackage{fontspec}
\setmainfont{Latin Modern Roman}
\usepackage{amsmath,amssymb,mathtools}
\usepackage{unicode-math}
\setmathfont{Latin Modern Math}
\usepackage{xurl,hyperref}
\hypersetup{colorlinks=true,urlcolor=blue}
\setcounter{secnumdepth}{0}
\setlength{\parindent}{0pt}
\setlength{\parskip}{5pt}
\setlength{\emergencystretch}{3em}
\newcommand{\sref}[2]{\hyperlink{source:#1:#2}{[S#2]}}
\newcommand{\cataloguescope}[1]{\paragraph{Recorded target.}#1}
\newcommand{\R}{\mathbb R}
\newcommand{\C}{\mathbb C}
\newcommand{\N}{\mathbb N}
\newcommand{\Z}{\mathbb Z}
\newcommand{\Q}{\mathbb Q}
\newcommand{\F}{\mathbb F}
\newcommand{\eps}{\varepsilon}
\begin{document}
\section{Variance Feedback in Cloud Boundary-Layer Models}\label{statement}

\subsection{Definitions and mathematical statement}

\paragraph{Shared notation.}$\mathbb N=\{1,2,\ldots\}$; the standard number fields have their usual meanings. All additional conventions are stated locally.


\paragraph{Definitions and conventions.}
A Reynolds decomposition writes a fluctuating quantity as its mean plus a fluctuation, and variance is the mean square of that fluctuation. Liquid-water potential temperature and total-water mixing ratio are the two cloud variables in the source closure equations. Their variance budgets contain production, transport and dissipation terms. A slow positive feedback is a coupling in which a change in a mean field enhances variance production, which in turn reinforces that change.

\paragraph{Target retained from the canonical source catalogue.}
\paragraph{Catalogue wording retained for traceability.}
For Reynolds-averaged equations modeling the variances of liquid-water potential temperature and total-water mixing ratio in a precipitating cloudy atmospheric boundary layer, can the equations represent explicitly the slow positive feedback that produces the variance without an additional stochastic submodel?

\paragraph{Source scope and qualification.}
The report\textquotesingle s precipitating boundary-layer budget and feedback mechanism are the target; the extracted statement does not supply the full closure equations. Adding an independent stochastic production model would not meet the stated without-an-additional-submodel requirement.
\subsection{Short English statement}
Can deterministic Reynolds-averaged cloud equations explicitly generate the observed slow feedback that raises thermodynamic variances?
\subsection{Research attempt}
\paragraph{Model, date and investigated scope.}
OpenAI GPT-6 Astra (configured model \texttt{gpt-6-astra}), ultra reasoning; generated 11 October 2026. Proved positivity and zero-covariance invariance for a deterministic homogeneous fluctuation closure, isolating the need for a physically derived production term.

\paragraph{Sources and literature check.}
Read the supplied definition, catalogue wording, source qualifications and reference entries for this question. The historical source/context and inherited status are retained; no fresh literature or novelty certification is claimed. No external theorem is used without proof in this attempt; explicit model assumptions are distinguished from facts about the source model.

\paragraph{Can deterministic feedback create variance from zero?} Test a two-variable fluctuation model \(\dot\xi=A(t)\xi\), where \(A\) is a continuous real \(2\times2\) matrix that may depend on deterministic mean fields. With mean-zero random initial fluctuations and covariance \(C=\mathbb E[\xi\xi^T]\), differentiation gives
\[
 \dot C=AC+CA^T,\qquad C(t)=\Phi(t)C(0)\Phi(t)^T,
\]
where \(\dot\Phi=A\Phi\), \(\Phi(0)=I\). The formula is verified by differentiation; uniqueness follows from the linear matrix equation. It preserves positive semidefiniteness because \(v^TC(t)v=(\Phi^Tv)^TC(0)(\Phi^Tv)\ge0\).
\paragraph{The zero-covariance obstruction.} In particular \(C(0)=0\) implies \(C(t)=0\) for every finite time for which the coefficients exist. Positive feedback through the means can amplify a pre-existing covariance but cannot seed it in this homogeneous closure. In the scalar constant case \(\dot v=2av\), the exact solution \(v=v_0e^{2at}\) makes both assertions explicit, including \(a>0\).
\paragraph{What is missing.} A deterministic production or transport term \(Q\) could break the invariance, but it must be derived from the source's precipitating boundary-layer budgets; calling an arbitrary chosen \(Q\) the physical mechanism would beg the question. A stochastic submodel is not necessary for the algebraic possibility, yet no suitable physical deterministic closure is constructed here. The first gap is identifying and closing the real temperature/water covariance production while maintaining positivity, conservation and the observed slow time scale. This diagnostic does not claim that the report's actual equations are homogeneous.

\paragraph{Verification and adversarial audit.}
Verified the covariance evolution, nonnegative matrix preservation and zero-variance case, while avoiding a claim that the actual cloud equations have this simplified form. This is a GPT-6 Astra Ultra self-audit, not independent human verification. Prompt A drafting followed by prompt B scope/dependency/boundary self-audit. Verified the covariance evolution, nonnegative matrix preservation and zero-variance case, while avoiding a claim that the actual cloud equations have this simplified form. The remaining gap is: Derive and close the source precipitation/transport variance production with the required slow feedback and positivity. Unresolved status and zero original-target completion are retained.

Coordinator checked the covariance evolution formula, positivity and invariance of zero covariance in the homogeneous deterministic closure. This additional check was model-assisted, not independent human verification.

\paragraph{Tools, process and contributions.}
Python 3 and local file-reading/writing tools for manuscript serialization and structural checks; mathematical calculations were symbolic reasoning, with no numerical experiment or theorem prover. Contribution prompt A was followed by prompt B. GPT-6 Astra with ultra reasoning drafted the research and performed its stated self-audit; any additional coordinator contribution is identified in the verification and audit. A GPT-6-based Codex coordinator assembled the manuscript and checked repository integration; its exact variant and reasoning configuration were not exposed. No human mathematical contribution, independent human proof check or formal proof certification is supplied. Complete derivations and the consulted source references are included here; private reasoning transcripts are not published.

\paragraph{Final status and first remaining gap.}
\textbf{UNRESOLVED}. Derive and close the source precipitation/transport variance production with the required slow feedback and positivity.

\section{Completion Estimate}
\noindent\textbf{COMPLETION: 0\%}
The proved calculation is an elementary auxiliary result or a test of a restricted strategy; no new portion of the full source target is established.

\subsection{Sources}
\begin{itemize}
\item[\hypertarget{source:statement:1}{S1}] ULAM.AI and the UnsolvedMath Contributors, UnsolvedMath, catalogue entry. Catalogue curation under CC BY 4.0; definitions and English exposition added locally.
\url{https://huggingface.co/datasets/ulamai/UnsolvedMath}.
\item[\hypertarget{source:statement:2}{S2}] Moist Processes in the Atmosphere (2019), Oberwolfach Reports, DOI 10.4171/owr/2019/7, pp. 465--512. Primary report accessed 2026-10-11; relevant PDF pages 14.
\url{https://doi.org/10.4171/owr/2019/7}.
\item[\hypertarget{source:statement:3}{S3}] Stochastic parameterization review (catalogue-listed context; not independently consulted in this preparation).
\url{https://journals.ametsoc.org/view/journals/bams/98/3/bams-d-15-00268.1.xml}.
\end{itemize}
\label{end:statement}
\end{document}
